\begin{problem}[第31届全国中学生物理竞赛复赛][哈勃定律与红移]{24}
天文观测表明，远处的星系均离我们而去。著名的哈勃定律指出，星系离开我们的速度大小 $v = HD$，其中 $D$ 为星系与我们之间的距离，该距离通常以百万秒差距（Mpc）为单位；$H$ 为哈勃常数，最新的测量结果为 $H = 67.80\mathrm{km/(s\cdot Mpc)}$。当星系离开我们远去时，它发出的光谱线的波长会变长（称为红移）。红移量 $z$ 被定义为 $z = \frac{\lambda' - \lambda}{\lambda}$，其中 $\lambda'$ 是我们观测到的星系中某恒星发出的谱线的波长，而 $\lambda$ 是实验室中测得的同种原子发出的相应的谱线的波长，该红移可用多普勒效应解释。绝大部分星系的红移量 $z$ 远小于1，即星系退行的速度远小于光速。在一次天文观测中发现从天鹰座的一个星系中射来的氢原子光谱中有两条谱线，它们的频率 $\nu'$ 分别为 $4.549\times 10^{14}\mathrm{Hz}$ 和 $6.141\times 10^{14}\mathrm{Hz}$。由于这两条谱线处于可见光频率区间，可假设它们属于氢原子的巴尔末系，即为由 $n > 2$ 的能级向 $k = 2$ 的能级跃迁而产生的光谱。（已知氢原子的基态能量 $E_0 = -13.60\mathrm{eV}$，真空中光速 $c = 2.998\times 10^{8}\mathrm{m/s}$，普朗克常量 $h = 6.626\times 10^{-34}\mathrm{J}\cdot\mathrm{s}$，电子电荷量 $e = 1.602\times 10^{-19}\mathrm{C}$）
    \begin{subproblem}
        该星系发出的光谱线对应于实验室中测出的氢原子的哪两条谱线？它们在实验室中的波长分别是多少？
    \end{subproblem}
    \begin{subproblem}
        求该星系发出的光谱线的红移量 z 和该星系远离我们的速度大小 v；
    \end{subproblem}
    \begin{subproblem}
        求该星系与我们的距离 $D$。
    \end{subproblem}
\end{problem}